% MT057 / explanation — canonical editable content.
这个目标叫Elastic Net：把L1的稀疏作用和L2的收缩作用放在一起。对单个权重\(w_j\)，惩罚为\(\alpha|w_j|+\beta w_j^2\)，通常取\(\alpha,\beta\ge0\)；讨论两种作用同时存在时，两者应为正。



\begin{itemize}

\item \textbf{A依原解不选。} “一视同仁”不够精确。各坐标可以用相同的\(\alpha,\beta\)，但大小不同的权重付出的代价不同；配套答案按后一个意思判断。

\item \textbf{B正确。} 平方项对大权重增加得更快。例如绝对值从1变2时，L1项翻倍，平方项却变成4倍；L2因此对大系数施加更强收缩。

\item \textbf{C一般不成立。} 非零L1项在0处有尖点，普通多维问题不能直接使用OLS式的单次矩阵求逆，需要适当迭代算法。特殊一维或正交设计可以有显式解，不应把原解扩成“任何情况都没有解析表达”。

\item \textbf{D正确。} L1能使一些系数恰好归零，完成特征选择；L2还能缩小保留下来的大系数，二者同时作用。

\item \textbf{E错误。} 绝对值和平方对正负同等大小的权重给相同惩罚，没有禁止负权重。负系数只是表示该特征增大时预测向下变化，具体取决于编码与其他条件。

\end{itemize}


特征尺度会改变惩罚的实际效果，通常在训练流程中先做合适标准化。


\textbf{出处：} 2023B期中 Q6，原卷第3页（5分）；2025A期中 Q6，原卷第3页（5分）


\textbf{答案依据：} 2023B期中Q6随卷答案BD；2025A Q6同题。A、C的文字边界在此说明。
